\section{Complete edge extraction at a finite physical count}
\label{sec:finite-count-extraction}
The individual critical-word calculation does not describe singular
itinerary boundaries. We now include all such boundaries in a finite-count
extraction. The argument integrates the full finite collision graph,
rather than extending an inverse-coarea Jacobian through grazing.
It gives an integrable Fourier remainder for every fixed count cutoff.
The constants of this remainder retain their dependence on the cutoff;
they will not be treated as uniform long-time estimates.

Throughout this section $R\in I$, $n\ge1$, and $L\ge1$ are fixed.
The section is the actual $Y_R^*$, the measure is its probability
$\nu_R^*=c_*^{-1}\nu|_{Y_R^*}$, and $J_{n,R}$ is unchanged. Write
$\ell=(k,m)\in\Z^2\times\mathbb N$ for displacement and physical
collision count. Define
\begin{equation}\label{eq:finite-count-measure}
 \mu_{n,R}^{w,L}=(J_{n,R})_*
       \bigl(w\one_{\{N_{n,R}\le L\}}\nu_R^*\bigr).
\end{equation}
The weight is not normalized. In particular $w=1$ gives the restricted
original law, not the law conditioned on the cutoff event.

\subsection{The whole finite collision graph}
\begin{definition}[Finite-record subanalytic weights]
\label{def:finite-record-weight}
A bounded complex weight is admissible at $(n,L,R)$ if its real and
imaginary parts, on every finite collision chart with $N_{n,R}\le L$,
are globally subanalytic functions of the initial coordinates. Values
on null collision singularities may be fixed arbitrarily. Denote its
supremum bound by $M_w$. This class includes the weight one.
\end{definition}
Here globally subanalytic has its usual real-geometric meaning; it is
not a regularity assertion about every bounded-variation function.
Finite Boolean combinations of subanalytic conditions on a finite list
of actual return states and partial records up to the $n$th return
belong to this class.
Indeed, after fixing the successive return counts, each state is a
specified finite collision iterate. There are only finitely many such
choices below $L$. Products of bounded subanalytic initial, intermediate
and terminal weights are allowed for the same reason. No assertion
about their variation uniformly in the number of observations is made.

\begin{lemma}[Finite graph and actual mixed density]
\label{lem:finite-graph-density}
The measure \eqref{eq:finite-count-measure} has a compactly supported
mixed density $p_{n,R}^{w,L}(\ell,t)$. It has only finitely many nonzero
lattice components. For every fixed $\ell$ and admissible $w$, the real
and imaginary parts of this density have constructible representatives
in the sense of Cluckers--Miller: finite sums of products of globally
subanalytic functions and logarithms of positive globally subanalytic
functions. Each representative is analytic off a finite set.
\end{lemma}
\begin{proof}
Use finitely many bounded rational angular charts and disjoint measurable
chart pieces. Uniform finite horizon gives a finite set of possible
next-disk labels. The finite physical graph in the proof of
Lemma~\ref{lem:finite-record-variation} specifies flight equations, reflection,
positive entry roots and the exclusion of an earlier collision with any
candidate disk. For a fixed word of at most $L$ flights, the graph is
semialgebraic in its initial and auxiliary variables. Projection gives
semialgebraic finite collision records. Return membership adds the exact
rectangle tests at every collision, not just the terminal test. Angular
intervals become finite rational-chart intervals; at fixed $R$ their
endpoints are fixed real coefficients. Thus the events of exactly $n$
returns, total count $m\le L$, and displacement $k$ are finite unions of
semialgebraic chart pieces. The collision probability in a rational
angular coordinate has a bounded semialgebraic density. In particular,
the construction includes the limiting grazing, competing-root and
return-membership boundaries through the domains of integration.

On a regular word, the total flight time is analytic.
Proposition~\ref{prop:general-edge} identifies its critical points as
normal-to-normal nondegenerate minima; in particular its gradient is
nonzero almost everywhere. That proposition is a physical collision-word
statement and does not depend on which measurable section is used to
restrict the initial state. Apply coarea on an exhaustion where the
gradient is bounded below, multiply by the bounded weight, and then sum
the finitely many words. This gives the density, with the normalization
$1/(4\pi c_*)$ in $(\alpha,p)$ coordinates. Singular initial states and
chart boundaries have zero initial measure. Restricting to $Y_R^*$
does not import the old section's return counts. The time support lies
in $[0,L\tau_{\max}]$; the count and displacement labels also range over
finite sets.

We identify the density without bounding an inverse Jacobian near a
singularity. On each chart let $b(x)$ be the product of its measure
density and $w$, and let $\Omega_\ell$ be its exact finite-record event.
For real $t$, the cumulative component is a sum of integrals
\begin{equation}\label{eq:constructible-cumulative}
 P_\ell(t)=\sum_{\rm charts}\int
   \one_{\Omega_\ell}(x)b(x)\one_{\{\tau_m(x)<t\}}\dd x .
\end{equation}
Extend the chart integrands by zero. They are globally subanalytic and
integrable for every $t$. The stability-under-integration theorem
\cite[Theorem 1.3]{CMInt} makes $P_\ell$ constructible. The absolute
continuity just proved identifies its almost-everywhere derivative with
$p_{n,R}^{w,L}(\ell,\cdot)$.

For completeness, only one-variable differentiation is required here.
Write a constructible representation as
$\sum_i a_i(t)\prod_j\log b_{ij}(t)$, with $b_{ij}>0$.
A common finite subanalytic partition makes all generators analytic on
each open interval. Their derivatives are subanalytic there, and the
product rule, with $(\log b)'=b'/b$, gives a constructible derivative.
Define its value arbitrarily at the finitely many endpoints. This is a
constructible representative of the density. The same argument treats
real and imaginary parts and proves the last assertion.
\end{proof}

\begin{remark}[Relation to the retained physical edge calculation]
\label{rem:finite-count-inherited-edge}
The regular critical points and their exponentially small individual
jumps were already classified in Proposition~\ref{prop:general-edge}.
They are not new consequences of the present lemma. The added conclusion
is constructible structure for the \emph{entire} finite-count component,
including all boundaries of its exact integration domain. A singular
boundary may produce a fractional power or a logarithm instead of a
finite jump. The extraction below allows all of these forms.
\end{remark}

\subsection{One-sided power--logarithm jets}
\begin{lemma}[A finite jet removes every second-derivative singularity]
\label{lem:power-log-jet}
Let $f\in L^1(\R;\mathbb C)$ have compact support and constructible
real and imaginary parts. There is a finite set $S\subset\R$ such that
$f$ is analytic off $S$. For every $s\in S$ and sign
$\epsilon\in\{-1,1\}$, on some interval $0<x<d_s$ one has
\begin{equation}\label{eq:power-log-expansion}
 f(s+\epsilon x)=\sum_{k=0}^{K_s}(\log x)^k
                         \sum_{j=j_s}^{\infty}c_{s,\epsilon,j,k}x^{j/q_s},
 \qquad q_s\in\mathbb N.
\end{equation}
The inner series are convergent Laurent series in $x^{1/q_s}$; their
negative parts are finite. After collecting equal exponents, every
nonzero coefficient with $j/q_s\le-1$ vanishes. Put
\begin{equation}\label{eq:power-log-jet}
 J_{s,\epsilon}(x)=
    \sum_{-1<j/q_s\le1}\sum_{k=0}^{K_s}
                           c_{s,\epsilon,j,k}x^{j/q_s}(\log x)^k.
\end{equation}
Then the difference in \eqref{eq:power-log-expansion} after subtracting
this finite jet has, for $a=0,1,2$, a bound
\begin{equation}\label{eq:power-log-remainder}
 \left|\frac{d^a}{dx^a}
       \bigl(f(s+\epsilon x)-J_{s,\epsilon}(x)\bigr)\right|
 \le C_{s,\epsilon}x^{\beta_s-a}(1+|\log x|^{K_s}),
 \qquad \beta_s>1,
\end{equation}
after decreasing $d_s$. An identically zero remainder satisfies the
same assertion with any $\beta_s>1$.
\end{lemma}
\begin{proof}
The one-variable specialization of constructible preparation
\cite[Theorem 3.11]{CMInt} writes a germ as a finite sum of rational
powers, nonnegative integral logarithmic powers, and subanalytic units.
At a one-sided endpoint, a strong subanalytic unit is analytic in
finitely many bounded rational monomials. Choose a common denominator
for these monomials and the rational prefactors. Expanding the analytic
units gives convergent Laurent series in the corresponding root of $x$.
Equivalently, one may first take the convergent one-variable Puiseux
germs of the subanalytic generators: if a positive generator is
$x^r u(x^{1/q})$ with $u(0)>0$, its logarithm is
$r\log x+\log u(x^{1/q})$. This gives exactly
\eqref{eq:power-log-expansion}, not merely an undifferentiated asymptotic
formula. Finitely many generators and partition endpoints suffice for
all of $f$. Include the endpoints of its support in $S$.

For the integrability assertion, take the smallest exponent having a
nonzero coefficient polynomial in $\log x$. Its highest nonzero log
power dominates all the lower log powers; every strictly higher
rational exponent is smaller by a positive power of $x$ times a fixed
logarithmic power. If that smallest exponent were at most $-1$, the
absolute value would not be integrable at zero. This contradicts
$f\in L^1$. Repeating the argument, or simply using minimality, excludes
all such exponents. Equal-exponent cancellations have already been
performed. The argument works also for complex coefficients.

Only finitely many terms have exponent at most one. After they are
removed, each of the finitely many log coefficients is either zero or
$x^{\beta}a(x^{1/q_s})$ with $\beta>1$ and $a$ analytic at zero.
Differentiate this representation twice. Bounded derivatives of $a$ on
a smaller root interval give \eqref{eq:power-log-remainder}, with the
least of the remaining exponents as $\beta_s$. In particular, termwise
differentiation is justified by the convergent analytic representation.
\end{proof}

Choose disjoint neighborhoods $(s-d_s,s+d_s)$, with $0<d_s<1$.
Let $\chi_s(x)$ be a $C^2$ piecewise-polynomial cutoff on $[0,\infty)$,
equal to one for $x\le d_s/2$ and zero for $x\ge d_s$. Its first two
derivatives vanish at the transition endpoints. Define
\begin{equation}\label{eq:full-finite-jet-density}
 e_f(t)=\sum_{s\in S}\sum_{\epsilon=\pm1}
  \one_{\{\epsilon(t-s)>0\}}\chi_s(\epsilon(t-s))
                         J_{s,\epsilon}(\epsilon(t-s)),
 \qquad r_f=f-e_f.
\end{equation}
The cutoffs need not be analytic; their piecewise-polynomial choice
keeps them subanalytic and is sufficient for two weak derivatives.

\begin{proposition}[Exact subtraction and its derivative budget]
\label{prop:finite-jet-subtraction}
One has $e_f\in L^1(\R)$ and $r_f\in W^{2,1}(\R)$, both with compact
support. The representative of $r_f$ is $C^1$ and has value and first
derivative zero at every point of $S$.

The budget for its second derivative can be recorded without suppressing
its endpoint constants. Put
\[
 \mathcal J(a,k,d)=\int_0^d x^{a-1}|\log x|^k\dd x
 =d^a\sum_{j=0}^k\frac{k!}{(k-j)!}
                      \frac{(-\log d)^{k-j}}{a^{j+1}}
 \quad(a>0,\ 0<d<1).
\]
On $0<x<d_s/2$, the contribution to $\|r_f''\|_1$ is bounded by
\begin{equation}\label{eq:explicit-jet-budget}
 C_{s,\epsilon}\left(
   \mathcal J(\beta_s-1,0,d_s/2)
  +\mathcal J(\beta_s-1,K_s,d_s/2)\right).
\end{equation}
On each transition interval the contribution is bounded by the integral
of $|f''|$ there plus
\[
 \sum_{j,k}|c_{s,\epsilon,j,k}|C_{j/q_s,k}\,
                d_s^{j/q_s-1}(1+|\log d_s|^k).
\]
On the remaining compact regular intervals it is the integral of
$|f''|$. Thus all constants in this finite budget are attached to actual
germs and actual regular intervals.
\end{proposition}
\begin{proof}
Every extracted exponent exceeds $-1$, so $e_f$ is integrable.
Near $s$, \eqref{eq:power-log-remainder} gives zero traces for both
$r_f$ and its first derivative on both sides, and gives an integrable
second derivative since $\beta_s-2>-1$. Away from these points all
functions are $C^2$, with compact support. The matching traces show by
integration by parts that no point mass or derivative of a point mass
remains in the second distributional derivative. This proves
$r_f\in W^{2,1}$ and the asserted representative.

Integrating \eqref{eq:power-log-remainder} for $a=2$ gives
\eqref{eq:explicit-jet-budget}. The identity for $\mathcal J$ follows
by repeated integration by parts, with the boundary term at zero equal
to zero. On $[d_s/2,d_s]$, two differentiations of
$\chi_s(x)x^\alpha(\log x)^k$ and the scaled cutoff bounds give the
stated transition estimate. The remaining regular intervals are
compact subsets of analytic pieces, so their displayed integrals are
finite. These are bounds for the cancellation remainder, not separate
integrals of the divergent second derivatives at $s$.
\end{proof}

\begin{remark}[Why constants and slopes are extracted]
\label{rem:constant-slope-extraction}
Even when the classical second derivative of a one-sided constant or
linear term vanishes, its zero extension has a distributional boundary
term. The constant produces a derivative of a point mass, and the slope
produces a point mass. Including both exponents $0$ and $1$ in
\eqref{eq:power-log-jet} gives an $L^1$ second derivative, stronger than
merely a finite-measure second derivative. Terms such as
$x^{1/2}$, $\log x$, or $x\log x$ must likewise be extracted. The
selection rule concerns the prepared raw density, not the variation of
an initial-coordinate observable.
\end{remark}

\subsection{All finite-count components at once}
\begin{theorem}[Complete finite-count raw extraction]
\label{thm:finite-count-raw-extraction}
For each fixed $n,L,R$ and every admissible weight, there is an exact
identity of finite complex measures
\begin{equation}\label{eq:finite-count-exact-extraction}
 \mu_{n,R}^{w,L}=E_{n,R}^{w,L}+Q_{n,R}^{w,L},
\end{equation}
where $E$ has the finite power--logarithm density
\eqref{eq:full-finite-jet-density} in each nonzero lattice component,
and the density $r_{n,R}^{w,L}(\ell,\cdot)$ of $Q$ belongs to
$W^{2,1}(\R)$. In particular
\begin{equation}\label{eq:finite-count-residual-norms}
 A_j(n,L,R,w)=\sum_\ell\|\partial_t^j
                              r_{n,R}^{w,L}(\ell,\cdot)\|_1<\infty
 \quad(j=0,2).
\end{equation}
Its mixed Fourier transform satisfies
\begin{align}
 \int_{\T^3\times\R}|\widehat Q|\dd\omega
 &\le4(2\pi)^3\sqrt{A_0A_2},
       \label{eq:finite-count-Fourier-integrable}\\
 \int_{\T^3\times\{|b|>B\}}|\widehat Q|\dd\omega
 &\le\frac{2(2\pi)^3}{B}A_2\qquad(B>0).
       \label{eq:finite-count-far-tail}
\end{align}
The construction includes every regular critical, singular itinerary
and image-boundary contribution to the restricted measure. It does
not assume that its only singularities are the retained Morse jumps.
\end{theorem}
\begin{proof}
Apply Lemma~\ref{lem:finite-graph-density} and
Proposition~\ref{prop:finite-jet-subtraction} to each of the finitely
many lattice components. Alternatively apply them to every chart-word
subdensity and then sum; that version records every word's derivative
budget separately. In either construction all actual finite graph
pieces are integrated in \eqref{eq:constructible-cumulative}. Boundary
points themselves have zero initial measure, but their effects on
one-sided density limits are included in these integrals. The prepared
germs of the whole density include those effects. There is no assumption
that an inverse coarea map extends smoothly across such a boundary.

For a component $r_\ell$, distributional integration by parts gives
\[
 |\widehat r_\ell(b)|\le
       \min\{\|r_\ell\|_1,|b|^{-2}\|r_\ell''\|_1\}.
\]
The integral on the real line of $\min\{a,d/b^2\}$ is
$4\sqrt{ad}$ for $a,d>0$. Sum over $\ell$, apply Cauchy--Schwarz,
and integrate the torus of volume $(2\pi)^3$ to obtain
\eqref{eq:finite-count-Fourier-integrable}. A zero norm is handled by
continuity, or by observing that a compactly supported affine $L^1$
function is zero. Integration only over $|b|>B$ gives the other bound.
\end{proof}

The constants in \eqref{eq:finite-count-residual-norms} are finite for
each fixed packet. Neither the preparation theorem nor compactness of
$I$ by itself bounds them uniformly in $R$ as $n,L$ grow. Small cutoff
widths, approaching singular values, jet coefficients and the number of
pieces all occur in Proposition~\ref{prop:finite-jet-subtraction}.
Their total contribution must be estimated before using
\eqref{eq:finite-count-far-tail} in a long-time frequency splice.
